% Part: computability
% Chapter: computability-theory
% Section: total

\documentclass[../../../include/open-logic-section]{subfiles}

\begin{document}

\olfileid{cmp}{thy}{tot}
\olsection{सर्वत्र परिभाषितता अनिर्णेय आहे}

एखादी गोष्ट असंगणनक्षम आहे हे दाखवण्यासाठी $s$-$m$-$n$
प्रमेय वापरण्याचे आणखी एक उदाहरण पाहू. $\fn{Tot}$ हा
सर्वत्र परिभाषित संगणनक्षम फलनांच्या निर्देशांकांचा संच असू
द्या, म्हणजे
\[
\fn{Tot} = \Setabs{x}{\text{प्रत्येक $y$ साठी, $\cfind{x}(y)\fdefined$}}.
\]

\begin{prop}
\ollabel{prop:total}
$\fn{Tot}$ संगणनक्षम नाही.
\end{prop}

\begin{proof}
$\fn{Tot}$ संगणनक्षम नाही हे पाहण्यासाठी $K$ चे त्याच्याकडे
न्यूनीकरण होते हे दाखवणे पुरेसे आहे. $h(x,y)$ ची व्याख्या
पुढीलप्रमाणे करा:
\[
h(x,y) \simeq
\begin{cases}
0 & \text{जर $x \in K$ असेल तर} \\
\fundefined & \text{अन्यथा}
\end{cases}
\]
$h(x,y)$ हे $y$ वर अजिबात अवलंबून नाही, हे लक्षात घ्या.
$h$ हे आंशिक संगणनक्षम आहे हे पाहणे कठीण नसावे: आदान
$x, y$ वर~$h$ संगणित करताना आधी~$\cfind{x}$ चे
आदान~$x$ वरील अनुकरण करतो; हे संगणन थांबले, तर
$h(x,y)$ हे $0$ परत करून थांबते. म्हणून $h(x,y)$ हे
$\Zero(\umin{s}{T(x,x,s)})$ इतकेच आहे; येथे $\Zero$
हे स्थिर शून्य फलन आहे.

$s$-$m$-$n$ प्रमेयानुसार असे आदिम पुनरावर्ती फलन~$k(x)$
अस्तित्वात आहे की प्रत्येक $x$ आणि~$y$ साठी
\[
\cfind{k(x)}(y) =
\begin{cases}
0 & \text{जर $x \in K$ असेल तर} \\
\fundefined & \text{अन्यथा}
\end{cases}
\]
म्हणून $x \in K$ असल्यास $\cfind{k(x)}$ सर्वत्र परिभाषित
असते; अन्यथा ते अपरिभाषित असते. अशा प्रकारे $k$ हे $K$
चे~$\fn{Tot}$ कडे न्यूनीकरण आहे.
\end{proof}

\begin{digress}
खरे तर $\fn{Tot}$ संगणनक्षमपणे प्रगणनीयही नाही; त्याची
गुंतागुंत ``अंकगणितीय पदानुक्रमात'' अधिक वरच्या स्तरावर
येते. मात्र येथे आपण या अधिक मजबूत निष्कर्षाचा विचार करणार नाही.
\end{digress}

\end{document}
